\section{Latency and physical scale}
\label{sec:latency}
\subsection{The relevant clock}
A remote project competes with what can be experienced before its returns
arrive. Let $T_j^{\rm lag}>0$ denote its objective delay and let
$\omega(u)>0$ be subjective years per objective year. With a constant
subjective discount rate $\delta\geq0$, the correct discount exponent is
\begin{equation}
 \mathcal D_j(t)=\delta\int_t^{t+T_j^{\rm lag}}\omega(u)\dd u,
 \qquad r_j(t)=r_{j0}\exp[-\mathcal D_j(t)],\quad
 S_j(t)=S_{j0}\exp[-\mathcal D_j(t)].
 \label{eq:subjective-delay}
\end{equation}
The simpler product $\delta\omega(t)T_j^{\rm lag}$ is exact for a clock
constant over the delay, and a lower bound for a non-decreasing clock.
A clock measured in processor cycles becomes a subjective clock only
under an additional account of how those cycles implement valued
experience. Neither high throughput nor parallel world creation alone
establishes this identification.

\begin{proposition}[Latency sufficient for the logarithmic corner]
\label{prop:latency}
Take $\eta=1$, $\ell(z)=\ln(1+z)$, fixed finite non-negative base returns
$r_{j0},S_{j0}$, $d_0\geq1$, and finitely many directions. Suppose
$a(t)\geq a_{\min}>0$ and $p_j(t)\geq p_{\min}>0$. With
$\mathcal D_{\min}(t)=\min_j\mathcal D_j(t)$, a sufficient condition for eventual strict
all-inward allocation is
\begin{equation}
 \mathcal D_{\min}(t)-\ln\{[1+B(t)][1+\ln(1+B(t))]\}\longrightarrow\infty.
 \label{eq:latency-growth}
\end{equation}
If $B(t)\geq B_{\min}>0$ eventually, the expression inside the logarithm
can instead be $\Gamma(B)=(1+B)\ln(1+B)$. For a non-decreasing clock
and $T_j^{\rm lag}\geq T_{\min}^{\rm lag}>0$, either condition follows
from its version with $\mathcal D_{\min}$ replaced by
$\delta T_{\min}^{\rm lag}\omega(t)$.
For $B(t)\leq\bar B$, a fixed sufficient threshold is
\begin{equation}
 \mathcal D_{\min}(t)>
 \max\left\{0,\max_j\ln\frac{(1+\bar B)(r_{j0}+d_0S_{j0})/a_{\min}
                   +\Gamma(\bar B)S_{j0}}{p_{\min}}\right\},
 \label{eq:latency-bounded}
\end{equation}
where a zero numerator is omitted.
\end{proposition}
\noindent\emph{Proof sketch.}
Insert the discounted coefficients into \cref{eq:exact}; the required
price is the undiscounted expression multiplied by $\exp(-\mathcal D_j)$.
The two terms are uniformly bounded by a constant times the bracket in
\cref{eq:latency-growth}. Complete bounds and counterexamples to weaker
assumptions are in \cref{app:latency-proof}.

Latency can therefore remove the constant-residual premise of the growth
obstruction. This is a price-adjusted comparison, not an assertion that
remote projects cease to produce physical resources. Increasing local
speed and reducing travel delay act in opposite directions. Discounting
in objective time replaces $\mathcal D_j$ by $\delta_{\rm obj}T_j^{\rm lag}$;
then raising $\omega$ has no effect. The choice of clock is a preference
assumption with an observable consequence for the model's allocation,
not a conclusion forced by relativity or thermodynamics.

Replication itself can invalidate a finite-return specification. Suppose
one descendant generation takes the same subjective delay $D>0$, has
expected reproduction factor $\mathcal R_X>0$, and each descendant yields
one unit of payoff at birth. The discounted sum from generations
$k\geq1$ is finite exactly when
\begin{equation}
 \mathcal R_X\exp(-D)<1,
 \qquad \sum_{k\geq1}\mathcal R_X^k\exp(-kD)
       =\frac{\mathcal R_X\exp(-D)}{1-\mathcal R_X\exp(-D)}.
 \label{eq:discounted-replication}
\end{equation}
This follows from the geometric-series criterion, with divergence also
at equality. Physical descendants can grow while discounted value stays
finite. Conversely, timing the first launch without accounting for its
future descendants can substantially understate the effective return.
The finite base $r_{j0}$ in \cref{prop:latency} must represent the whole
specified payoff stream, or a genuinely finite planning horizon.

\subsection{Depletion and resilience are different mechanisms}
\begin{proposition}[No endogenous price drift on a stationary frontier]
\label{prop:depletion}
Suppose $p_j(t)=P_j(X(t))$ depends only on cumulative past external
investment, with $X(t)=X(t_0)$ along an all-inward path starting at $t_0$.
Then $p_j(t)=P_j(X(t_0))$ throughout that path. Frontier depletion alone
cannot supply a growing relative price along it.
\end{proposition}
\begin{proof}
Substitute the constant argument into each $P_j$. Historical expansion
can determine the level at $t_0$, but cannot create subsequent drift
without further investment or another state variable.
\end{proof}

\begin{corollary}[A conditional risk cost of discount-induced abstention]
\label{cor:quiet-risk}
Suppose the only available mitigation of a catastrophe has physical
success probability $\varphi(b)>0$ for $b>0$, baseline catastrophe
probability $\vartheta>0$, and no effect before its arrival. An allocation
with no such backup retains catastrophe probability $\vartheta$ after
that arrival date, while a completed backup $b$ reduces it to
$\vartheta[1-\varphi(b)]$. A utility discount can change which allocation
is preferred but does not reduce either physical probability.
\end{corollary}
\begin{proof}
These probabilities follow from the stipulated event tree. Multiplying
the valuation of a future outcome by a discount factor does not change
the probability assigned to that outcome.
\end{proof}
This is the defensible content of a fragility trade-off. There is no
general implication from quietness to fragility when dormant backups,
local mitigation, negligible hazards, or information about future hazards
are available. Only geographically mitigable, otherwise unaddressed risk
belongs in this comparison.

\paragraph{Remark (A finite-lived host).}
If the planning horizon includes certain terminal loss of the usable host,
and only geographical separation mitigates it, the stay-put risk model has
$\vartheta=1$. For positive computation and a feasible backup with
$\varphi'(0)>0$, \cref{app:risk-proof} excludes the inward corner.
Efficiency can postpone resource insufficiency, not remove this stipulated
terminal loss; discounting does not lower its probability
(\cref{cor:quiet-risk}).

\subsection{An explicit dimensional dictionary}
The return relevant to delay is a \emph{dimensionless marginal advantage},
not the bare normalized coefficient $r_j/p_j$:
\begin{equation}
 \mathcal R_{{\rm marg},j}
   =\frac{r_{j0}}{p_j a\ell'(B)},\qquad
 \mathcal D_j\gtrsim\ln\mathcal R_{{\rm marg},j}
 \label{eq:physical-return-ratio}
\end{equation}
for a pure-computation comparison with delayed external output. Residual
motives require the full preceding inequalities. Resource acquisition by
self-replicating industry can make large returns conceivable
\citep{ArmstrongSandberg2013}; that work does not provide a calibrated
value of this paper's normalized $r_j/p_j$. A ratio $10^{20}$ below is a
scenario assumption, not an empirical estimate attributed to that source.

For an energy-limited illustration, choose $b_{\rm ref}=10^{32}\,{\rm J}$
and a planning horizon of one Julian year, $31\,557\,600\,{\rm s}$.
Using the IAU nominal solar luminosity
$L_\odot^{\rm N}=3.828\times10^{26}\,{\rm W}$ \citep{IAU2015}, a captured
fraction $f_{\rm cap}$ gives
\begin{equation}
 \bar B=120.8024928\,f_{\rm cap}
 \left(\frac{L_*}{L_\odot^{\rm N}}\right)
 \left(\frac{\Delta}{1\,{\rm yr}}\right)
 \left(\frac{10^{32}\,{\rm J}}{b_{\rm ref}}\right).
 \label{eq:stellar-budget}
\end{equation}
This is a per-horizon flow, not the star's lifetime energy. The scalar
budget does not replace independent material, memory, or heat constraints.
A fixed ceiling describes a specified era and accessible substrate;
stellar evolution need not leave it constant indefinitely.

A distance of four light-years traversed at $0.1c$ has a forty-year
travel time by definition of a light-year. This is an illustrative
one-way journey, not a complete replication cycle; manufacture,
deceleration, and return information add delays. For
$\mathcal R_{\rm marg}=10^{20}$, $\omega=10^6$, and a forty-year delay,
\cref{eq:physical-return-ratio} requires
$\delta\gtrsim1.1513\times10^{-6}$ per subjective year.
Neither the assumed clock speed nor the discount rate is measured for
an extraterrestrial civilization. Their separate roles remain visible
in \cref{tab:latency-sensitivity} rather than being buried in a normalized
speed parameter.
\begin{table}[tbp]
\centering
\caption{Assumed marginal advantages and clock speeds: threshold
$\delta=\ln\mathcal R_{\rm marg}/(40\omega)$ per subjective year.
The forty-year delay and every return/clock entry are scenarios.}
\label{tab:latency-sensitivity}
\begin{tabular}{@{}rrrr@{}}\toprule
$\mathcal R_{\rm marg}$ & $\omega=1$ & $\omega=10^3$ & $\omega=10^6$\\\midrule
$10^3$ & $0.17269$ & $1.7269\times10^{-4}$ & $1.7269\times10^{-7}$\\
$10^{10}$ & $0.57565$ & $5.7565\times10^{-4}$ & $5.7565\times10^{-7}$\\
$10^{20}$ & $1.15129$ & $1.1513\times10^{-3}$ & $1.1513\times10^{-6}$\\
$10^{40}$ & $2.30259$ & $2.3026\times10^{-3}$ & $2.3026\times10^{-6}$\\\bottomrule
\end{tabular}
\end{table}

To interpret \cref{tab:ranges} physically, specify $b_{\rm ref}$,
$c_{\rm ref}$, a project-capacity unit for each direction, and the welfare
normalization. Then $B_{\rm res}=Bb_{\rm ref}$,
$C_{\rm useful}=Cc_{\rm ref}$, a unit project costs $p_jb_{\rm ref}$,
and $r_jc_{\rm ref}$ is its horizon computation. The local marginal
product is $a\ell'(B)c_{\rm ref}/b_{\rm ref}$. Benefit amplitudes and
scales require an independent valuation model; astronomical dimensions
do not identify them. Thus the original normalized interval $B\in[1,100]$
can represent $10^{32}$--$10^{34}\,{\rm J}$ under this dictionary, but
it is not intrinsically that energy range. Changing the dictionary
changes the interpretation, not the frozen experiment.
